\section*{第 \ref{ch:g9:fractions} 章——分数与幂}

\begin{solution}{exo:g9:fractions:1}
$\dfrac13 + \dfrac15 = \dfrac{5}{15} + \dfrac{3}{15} = \dfrac{8}{15}$.

$\dfrac72 - \dfrac53 = \dfrac{21}{6} - \dfrac{10}{6} = \dfrac{11}{6}$.

$\dfrac{5}{12} + \dfrac38 = \dfrac{10}{24} + \dfrac{9}{24} =
\dfrac{19}{24}$.

$3 - \dfrac45 = \dfrac{15}{5} - \dfrac45 = \dfrac{11}{5}$.
\end{solution}

\begin{solution}{exo:g9:fractions:2}
$\dfrac59 \times \dfrac{27}{35} = \dfrac{5 \times 27}{9 \times 35}
= \dfrac{27}{9} \times \dfrac{5}{35} = 3 \times \dfrac17 = \dfrac37$
（分别约去 $9$ 与 $5$）。

$\dfrac{8}{15} \div \dfrac45 = \dfrac{8}{15} \times \dfrac54
= \dfrac{8 \times 5}{15 \times 4} = \dfrac{40}{60} = \dfrac23$.

$\dfrac23 \times \dfrac34 \times \dfrac45 = \dfrac{2 \times 3 \times 4}
{3 \times 4 \times 5} = \dfrac25$（约去 $3$ 与 $4$）。
\end{solution}

\begin{solution}{exo:g9:fractions:3}
$2^5 = 32$; \quad $(-2)^4 = 16$; \quad $-2^4 = -16$; \quad
$4^{-2} = \frac{1}{16}$; \quad
$\left(\frac23\right)^3 = \frac{8}{27}$; \quad $7^0 = 1$.
\end{solution}

\begin{solution}{exo:g9:fractions:4}
$5^3 \times 5^6 = 5^9$; \quad $\dfrac{7^9}{7^5} = 7^4$; \quad
$\left(2^3\right)^4 = 2^{12}$; \quad
$\dfrac{3^2 \times 3^7}{3^5} = 3^{9-5} = 3^4$; \quad
$2^5 \times 4 = 2^5 \times 2^2 = 2^7$.
\end{solution}

\begin{solution}{exo:g9:fractions:5}
$45\,000 = 4.5 \times 10^4$; \quad
$0.0072 = 7.2 \times 10^{-3}$; \quad
$302.5 = 3.025 \times 10^2$; \quad
$0.000\,000\,9 = 9 \times 10^{-7}$; \quad
一千二百万 $= 1.2 \times 10^7$。
\end{solution}

\begin{solution}{exo:g9:fractions:6}
$(2 \times 10^7) \times (4.5 \times 10^{-3}) = 9 \times 10^{4}$.

$\dfrac{9 \times 10^5}{3 \times 10^{-2}} = 3 \times 10^{5 - (-2)}
= 3 \times 10^7$.

$(5 \times 10^3)^2 = 25 \times 10^6 = 2.5 \times 10^7$.
\end{solution}

\begin{solution}{exo:g9:fractions:7}
加入的水占水箱的
$\dfrac34 - \dfrac35 = \dfrac{15}{20} - \dfrac{12}{20} = \dfrac{3}{20}$
\omterm{def:g2:measure:units}{容积}。因此 $\frac{3}{20}$ 的\omterm{def:g2:measure:units}{容积}为 $30$ 升：
$\frac{1}{20}$ 对应 $10$ 升，\omterm{def:g2:measure:units}{容积}为 $200$ 升。
\end{solution}

\begin{solution}{exo:g9:fractions:8}
爱丽丝吃后，剩下蛋糕的 $1 - \frac14 = \frac34$。本
吃掉 $\frac13 \times \frac34 = \frac14$，剩下
$\frac34 - \frac14 = \frac12$。卡拉吃掉其中的\omterm{def:g3:division:half}{一半}，即 $\frac14$，
剩下蛋糕的 $\frac14$。
\end{solution}

\begin{solution}{exo:g9:fractions:9}
时间 $= \dfrac{\text{距离}}{\text{速度}} =
\dfrac{1.5 \times 10^{11}}{3 \times 10^8} = 0.5 \times 10^3 = 500$
秒，即 $8$ 分钟 $20$ 秒。
\end{solution}

\begin{solution}{exo:g9:fractions:10}
将两数都写成\omterm{def:g8:powers:def}{指数}为 $10$ 的幂：
\[
2^{30} = \left(2^3\right)^{10} = 8^{10},
\qquad
3^{20} = \left(3^2\right)^{10} = 9^{10}.
\]
由于 $8 < 9$，取十次幂得 $8^{10} < 9^{10}$，因此
$2^{30} < 3^{20}$.
\end{solution}

\begin{solution}{pb:g9:fractions:1}
\textbf{1.} $\frac14 = 0.25$（有限停止）；$\frac13 = 0.333\ldots$
（循环节为 $3$）；$\frac56 = 0.8333\ldots$（循环节为 $3$，在
$8$ 之后）；$\frac{3}{11} = 0.272727\ldots$（循环节为 $27$）。

\textbf{2.} 在以 $b$ 为除数的每一步中，\omterm{def:g3:division:remainder}{余数}
是 $b$ 个可能值 $0, 1, \dots, b - 1$ 中的一个
（见\cref{thm:g6:wholes:division}）。若\omterm{def:g3:division:remainder}{余数} $0$ 出现，则
除法停止。否则，至多经过 $b$ 步，
必有某个\omterm{def:g3:division:remainder}{余数}第二次出现——步数多于可能的\omterm{def:g3:division:remainder}{余数}个数——
从那时起，除法
恰好重复它在某\omterm{def:g3:division:remainder}{余数}首次出现之后所产生的数字序列，
即\omterm{def:g6:decimals:places}{小数}的书写无限循环下去。要么停止，
要么循环：不存在第三种情形。

\textbf{3.} $\frac17 = 0.142857\,142857\ldots$：循环节
$142857$ 的长度为 $6$。第 2 题保证在至多
$7$ 步之内出现重复——此处的\omterm{def:g3:division:remainder}{余数}
$1, 3, 2, 6, 4, 5$ 都在 $1$ 再次出现之前出现，
这是除以 $7$ 时可能的最长周期。

\textbf{4.} $\frac ab$ 停止，当且仅当某个
$\frac{a \times k}{b \times k}$ 的\omterm{def:g4:fractions:def}{分母}为 $10$、$100$、
$1000, \dots$——即\omterm{def:g4:fractions:def}{分母}能乘成
十的幂，如 $4 \times 25 = 100$ 或
$8 \times 125 = 1\,000$。对 $3$ 与 $6$ 而言这是行不通的：一个
十的幂的数字和为 $1$，而 $3$ 的每个\omterm{def:g5:division:multiple}{倍数}的数字和
都是 $3$ 的\omterm{def:g5:division:multiple}{倍数}——因此 $3$ 的\omterm{def:g5:division:multiple}{倍数}
（从而 $6$ 的\omterm{def:g5:division:multiple}{倍数}也同样）绝不可能是 $10$、$100$、$1\,000, \dots$。对于
$11$：$10$、$100$、$1\,000$、$10\,000$ 除以 $11$ 时总是
留下\omterm{def:g3:division:remainder}{余数}（依次为 $10$、$1$、$10$、$1, \dots$），从不为 $0$。
\omterm{def:g4:fractions:def}{分母}能乘成十的幂的\omterm{def:g9:fractions:fraction}{分数}便会停止；其余的都会循环。

\textbf{5.} $\frac{9}{40}$：$40 \times 25 = 1\,000$，停止：
$\frac{225}{1000} = 0.225$。$\frac{33}{50}$：
$50 \times 2 = 100$，停止：$\frac{66}{100} = 0.66$。
$\frac{5}{12}$：$12$ 是 $3$ 的\omterm{def:g5:division:multiple}{倍数}，循环
（$0.41666\ldots$）。

\textbf{6.} $10x = 7.777\ldots$，于是 $10x - x = 7$（尾部
恰好抵消），$9x = 7$，$x = \frac79$。验算除法：
$7 \div 9 = 0.777\ldots$。

\textbf{7.} 循环节有两位数字，故乘以 $100$ 移位：
$100x - x = 72.7272\ldots - 0.7272\ldots = 72$，因此
$99x = 72$，$x = \frac{72}{99} = \frac{8}{11}$。

\textbf{8.} $1000x = 583.333\ldots$，$100x = 58.333\ldots$；
相减得 $900x = 525$，于是
$x = \frac{525}{900} = \frac{7}{12}$（约去 $75$）。事实上
$\frac{7}{12} = 0.58333\ldots$。

\textbf{9.} $10x - x = 9.999\ldots - 0.999\ldots = 9$，于是
$9x = 9$，$x = 1$。并非“恰好小于 $1$”：
$0.999\ldots$ \emph{就是}数 $1$，只是换了一种
写法而已。任何严格小于 $1$ 的数都会留下空隙，而
由\cref{pb:g6:decimals:1}可知，空隙中总还含有其他
数——而 $0.999\ldots$ 与 $1$ 之间却容不下任何数。
无穷的九，正是
\cref{pb:g6:fractions:1}中那块巧克力缩小到比任何正数都小：
什么也没有剩下。

\textbf{10.} $2.4545\ldots = 2 + \frac{45}{99} = 2 +
\frac{5}{11} = \frac{27}{11}$。另外
$0.123123\ldots = \frac{123}{999} = \frac{41}{333}$（约去
$3$）。

\textbf{11.} 对照表：$0.037037\ldots = \frac{37}{999}$。
由于 $999 = 27 \times 37$，约分得 $\frac{1}{27}$：
\omterm{def:g9:fractions:fraction}{分数} $\frac{1}{27}$ 吟唱着“$037$”。验算除法：
$1 \div 27 = 0.037037\ldots$。

\textbf{12.} $0.0272727\ldots = \frac{1}{10} \times
0.272727\ldots = \frac{1}{10} \times \frac{27}{99}
= \frac{27}{990} = \frac{3}{110}$（约去 $9$）。

\textbf{13.} 零组成的数段无限增长，因此没有哪个固定的
数段能永远重复：它\emph{不是}循环\omterm{def:g6:decimals:places}{小数}
（它也永不停止）。由第一部分可知，每个\omterm{def:g9:fractions:fraction}{分数}都停止或
循环——故此数根本不是\omterm{def:g9:fractions:fraction}{分数}：它是
无理数，专门“定制”而成。无理性的超级明星，
$\sqrt2$，在\cref{ch:g9:sqrt}中被揭开真面目。

\textbf{14.} $2^{30} = \left(2^{10}\right)^3 \approx
(1.024)^3 \times 10^9 \approx 1.07 \times 10^9$：一个
略大于 $10^9$ 的数，因而有 $10$ 位数字。（精确地说：
$2^{30} = 1\,073\,741\,824$。）

\textbf{15.} $0.20262026\ldots = \frac{2026}{9999}$，已是
最简\omterm{def:g9:fractions:fraction}{分数}（$9999 = 3 \times 3 \times 11 \times 101$ 与
$2026 = 2 \times 1013$ 没有公\omterm{def:g5:division:multiple}{因数}）。完整的
对照表：\emph{\omterm{def:g9:fractions:fraction}{分数}恰好就是那些停止或
循环的\omterm{def:g6:decimals:places}{小数}}——\omterm{def:g4:fractions:def}{分母}能乘成十的幂者停止，其余
都吟唱一个循环节；反之，每一段“吟唱”，哪怕是你的
毕业年份，都是伪装的\omterm{def:g9:fractions:fraction}{分数}。
\end{solution}
